\section{Results}

Some results here.

Link to user guide (publish user guide on zenodo, link to actual dataset on zenodo) as reference for people who want details of how to use, not key results

Note key trends over various groups of gases

Comparison with CMIP6 and CMIP5 datasets

Compare to station data we started with (conclusion should be mostly that we are close in ERF terms, but there are clear areas for improvement or reconsideration)

Compare to any independent networks we can think of/find

Add discussion section

Note various issues identified in methods

Note issues for minor gases and lack of new datasets compared to last time, suggest picking a smaller, more clearly supported, set of gases for future to simplify things

Note that we still don't use satellite data: clear area for improvement

Note that we also don't include uncertainties. Note that leading uncertainty is likely not measurement-level,
rather methodological i.e. the various leaps/assumptions we have to make to go from a sparse observational network to a temporally extended, spatially complete dataset

Notes re extensions for later internal use:

- ingest updated network data
- calculate new global- annual-means
- optimise existing latitudinal gradient (and seasonality change for CO_2) to get best fit with interpolated observational network values
    - avoids the jump of using a different lat. gradient or seasonality over time
    - also avoids extending PCs using crude regressions
    - have to use relative seasonality calculation for gases other than CO_2, oh well just one less degree of freedom in the optimisation
- calculate full grids
- check no unexpected/spurious jumps
- process, format etc. then publish
